\section{Introduction}
Test-time compute has emerged as a key driver of progress in LLMs, with techniques like chain-of-thought reasoning and iterative refinement demonstrating that inference-time scaling can unlock new capabilities~\citep{wu2025inferencescalinglawsempirical,snell2024scalingllmtesttimecompute}.
The rapid rise of parallel, agentic workflows has only intensified the need for efficient inference and deployment of such models~\citep{openai_gpt53_codex_2026,anthropic_claude_opus_46_2026}.
This paradigm shift makes inference efficiency~\citep{kwon2023efficientmemorymanagementlarge,li2024llminferenceservingsurvey} paramount, as the practical impact of AI systems now depends critically on their ability to perform large-scale inference during deployment.
Model architecture design plays a fundamental role in determining inference efficiency, as architectural choices directly dictate the computational and memory requirements during generation.
While Transformer-based models~\citep{vaswani2017attention} are the current industry standard, they are fundamentally bottlenecked by linearly increasing memory demands through the KV cache and quadratically increasing compute requirements through the self-attention mechanism.
These drawbacks have motivated recent lines of work on sub-quadratic models, e.g., state space models (SSMs) and linear attention, which retain constant memory and linear compute while attaining comparable or better performance than their Transformer counterparts.
These models have made it into the mainstream, with layers such as Mamba-2~\citep{dao2024transformersssmsgeneralizedmodels} and Gated DeltaNet (GDN)~\citep{schlag2021deltarule,yang2025parallelizinglineartransformersdelta} recently incorporated into large-scale hybrid models that match the performance of pure Transformer alternatives with much higher efficiency~\citep{nvidia2025nemotronhfamilyaccurateefficient,qwen3technicalreport,kimiteam2025kimilinearexpressiveefficient,tencenthunyuanteam2025hunyuanturbosadvancinglargelanguage}.

Despite the success of linear models, significant progress remains in improving their performance, in particular on advancing the Pareto frontier between model quality and inference efficiency.
For example, Mamba-2 was developed to improve training speed and simplicity over Mamba-1~\citep{gu2024mambalineartimesequencemodeling}, by sacrificing some expressivity and thus performing worse for inference-matched models.
In addition, they have been shown to lack certain capabilities, such as poor state-tracking abilities, e.g., simply determining parity of bit sequences~\citep{grazzi2025unlockingstatetrackinglinearrnns,sarrof2024expressivecapacitystatespace}.
Finally, despite these sub-quadratic models being prized for theoretically efficient inference and thus their widespread adoption, their inference algorithms are not hardware efficient.
In particular, because these algorithms were developed from a training perspective, their decoding phase has low arithmetic intensity (the ratio of FLOPs to memory traffic), resulting in large portions of hardware remaining idle. 

To develop more performant models from an inference-first paradigm, we introduce three core methodological changes on top of Mamba-2, influenced by an SSM-centric viewpoint of sub-quadratic models.

\paragraph{Exponential-Trapezoidal Discretization.} 
We provide a simple technique for discretizing time-varying, selective SSMs.
Through our framework, we can derive several new discretization methods.
One of our instantiations, referred to as ``exponential-Euler,'' formalizes Mamba-1 and Mamba-2's heuristic discretization that previously lacked theoretical justification.
Our new ``exponential-trapezoidal'' instantiation is a more expressive generalization of ``exponential-Euler,'' where the recurrence can be expanded to reveal an implicit convolution applied on the SSM input. 
Combined with explicit $B,C$ bias terms,
Mamba-3 can empirically replace the short causal convolution in language model architectures, which was previously hypothesized to be essential for recurrent models.

\paragraph{Complex-valued State Space Model.} By viewing the underlying SSM of Mamba-3 as complex-valued, we enable a more expressive state update than Mamba-2's. This change in update rule, designed to be lightweight for training and inference, overcomes the lack of state-tracking ability common in many current linear models. We show that our complex-valued update rule is equivalent to a data-dependent rotary embedding and can be efficiently computed~\citep{su2023roformerenhancedtransformerrotary}, and empirically demonstrate its ability to solve synthetic tasks outside the capabilities of prior linear models.

\paragraph{Multi-Input, Multi-Output (MIMO) SSM.} To improve FLOP efficiency during decoding, we switch from an outer-product–based state update to a matrix-multiplication–based state update.
From the view of the signal processing foundations of SSMs, such a transition exactly coincides with the generalization from a single-input single-output (SISO) sequence dynamics to a multiple-input multiple-output (MIMO) one. Here, we find that MIMO is particularly suitable for inference, as the extra expressivity enables more computation during the memory-bound state update during decoding, without increasing the state size and compromising speed.

Put together, these improvements form the core of our \textbf{Mamba-3} layer.
Methodologically, we note that these all arise naturally from an SSM-centric perspective but are not immediate from other popular viewpoints of modern linear layers such as linear attention or test-time regression;
we discuss these connections further in~\Cref{sec:related}.
Empirically, we validate our new model's abilities and capabilities on a suite of synthetic state-tracking and language-modeling tasks.

\begin{itemize}[itemsep=0.1pt,topsep=0pt,leftmargin=*]
    \item \textbf{Better Quality.} 
    At 1.5B scale, Mamba-3 (MIMO) improves downstream language modeling accuracy by \textbf{+2.2} over Transformers, \textbf{+1.9 points} over Mamba-2, and \textbf{+1.8} over GDN, while Mamba-3 (SISO) improves over the next best model, GDN, by \textbf{+0.6} points.
    Furthermore, across state size experiments, Mamba-3 (MIMO) with state size 64 matches the perplexity of Mamba-2 with state size 128, effectively achieving the \textbf{same language modeling performance with half the latency}.
    \item \textbf{New Capabilities.} Mamba-3's complexification of the SSM state enables it to \textbf{solve synthetic state-tracking tasks that Mamba-2 cannot}. We empirically demonstrate that the efficient RoPE-like calculation is able to near perfectly solve arithmetic tasks, while Mamba-3 without RoPE and Mamba-2 perform no better than random guessing.
    \item \textbf{Inference Efficiency.} Mamba-3 (MIMO) improves hardware utilization. It increases decoding FLOPs by up to \textbf{4$\times$} relative to Mamba-2 at fixed state size, while maintaining \textbf{similar wall-clock decode latency}, and simultaneously improving perplexity and downstream performance. We release fast training and inference kernels for Mamba-3.\footnote{\url{https://github.com/state-spaces/mamba}.}
\end{itemize}

Mamba-3 (SISO) improves quality and capability over prior linear models, and Mamba-3 (MIMO) further improves performance over Mamba-3 (SISO) and other strong baselines while matching inference speed with Mamba-2.
Both of our Mamba-3 variants advance the performance-latency Pareto frontier through their strong modeling capabilities and hardware-efficient design.
